% ============================================================
% f2014_tianjin_arts_q07_circle.tex
% 2014 天津卷（文科）第 7 题：圆内接三角形与切线
% 来源 PDF：普通高考/2014/2014天津文.pdf
% 原页：第 1 页（A3），印面 925，第 7 题题干下方
% D 为 ∠BAC 平分线与圆的交点，E=AD∩BC，F 为过 B 的切线与 AD 延长线的交点。
% 取 ∠A=80°，B、C 的弧上位置 208°、338°，此配置下结论 ①②④ 成立、③ 不成立。
% ============================================================
\begin{tikzpicture}[>=Stealth, line width=0.7pt, font=\small,
  line cap=round, line join=round, scale=1.2]
  % 圆（半径 1）
  \draw[line width=0.7pt] (0,0) circle (1);
  \coordinate (A) at (0.174,0.985);
  \coordinate (B) at (-0.883,-0.469);
  \coordinate (C) at (0.927,-0.375);
  \coordinate (D) at (0.052,-0.999);
  \coordinate (E) at (0.088,-0.419);
  \coordinate (F) at (-0.015,-2.102);

  % 三角形与内接弦
  \draw (A) -- (B) -- (C) -- cycle;
  \draw (B) -- (D) -- (C);
  % 角平分线 AD 及延长线到 F
  \draw (A) -- (D);
  \draw (D) -- (F);
  % 过 B 的切线
  \draw (B) -- (F);

  \fill (A) circle (0.9pt) node[above right=-1pt] {$A$};
  \fill (B) circle (0.9pt) node[below left=-1pt] {$B$};
  \fill (C) circle (0.9pt) node[right=-1pt] {$C$};
  \fill (D) circle (0.9pt) node[below right=-1pt] {$D$};
  \fill (E) circle (0.9pt) node[above left=-1pt] {$E$};
  \fill (F) circle (0.9pt) node[below=-1pt] {$F$};
\end{tikzpicture}
